% !PRINTED 781
\weberpage{781}
%<<CONT
lich ist. Ihr Werth ist gleich der Quadratwurzel aus dem
Quotienten der Discriminanten
\[
\Delta (1, r, r^2, \ldots, r^{\nu-1}) : \Delta (\omega_1, \omega_2, \ldots, \omega_\nu),
\]
worin $\omega_1, \omega_2, \ldots, \omega_\nu$ eine Minimalbasis ist (\S.\,144, 145).

Wir schliessen alle in der Discriminante der Kreistheilungs\-gleichung und alle in $m$ und in $a$ aufgehenden Primzahlen, deren Anzahl jedenfalls endlich ist, von dem Systeme der $p$ aus, und bezeichnen mit $\mathfrak{p}$ ein in $p$ aufgehendes Primideal. Dann ist die Congruenz
\[
r^p \equiv r \ (\mathrm{mod}\ \mathfrak{p})
\]
immer dann, aber auch nur dann, erfüllt, wenn

\weq{16}{p \equiv 1 \ (\mathrm{mod}\ m)}

ist, d.\,h. wenn $p$ zum Exponenten 1 gehört, wenn also $r^p = r$ ist, und unter derselben Bedingung ist daher auch nach (15) für jede ganze Zahl $\omega$ in $\Omega_m$:
\[
\omega^p \equiv \omega \ (\mathrm{mod}\ \mathfrak{p}).
\]
Dies ist aber nach \S.\,162 die nothwendige und hinreichende Bedingung dafür, dass die in der Körperdiscriminante von $\Omega_m$ nicht aufgehende Primzahl $p$ in lauter Primideale ersten Grades zerfällt.

Bezeichnen wir also mit
\[
P_f(s) = \prod (1 - p^{-sf})
\]
das über alle zum Exponenten $f$ gehörigen Primzahlen $p$ erstreckte Product, so ergiebt sich aus \S.\,192, dass, wenn $f>1$ ist, $P_f(1)$ endlich und von Null verschieden ist, und dass

\weq{17}{\dfrac{s-1}{P_1(s)^{\nu}}}

\gesperrt{für} $s=1$ \gesperrt{endlich und von Null verschieden ist.}

Nun ist nach (13), wenn mit $S$ das über die \gesperrt{ausgeschlossenen} Primzahlen $p$ erstreckte \gesperrt{endliche} Product $\Pi(1-p^{-sf})^{-e}$ bezeichnet wird,
\[
Q_1 Q_2 \ldots Q_\mu = S \overset{f}{\Pi} \frac{1}{P_f(s)^{e}},
\]
und da $e = \mu$ für $f = 1$, so ergiebt sich aus (17):

\weq{18}{(s-1)^{\frac{\mu}{\nu}} Q_1 Q_2 \ldots Q_\mu}

\gesperrt{für} $s = 1$ \gesperrt{endlich und von Null verschieden.}
